\subsection{Continuous functions on a compact subset of C}

Lemma 1. --- Let X be a compact subset of the plane and let O be a bounded connected component of \(\mathbf{C} - \mathbf{X}\) . The boundary of O is contained in X.

The closure of the set O in C-X is equal to \(\overline{\mathrm{O}}\cap (\mathbf{C} - \mathbf{X})\) where \(\overline{\mathrm{O}}\) is its closure in C. Since O is a connected component of C-X, we therefore have \(\overline{\mathrm{O}}\cap (\mathbf{C} - \mathbf{X}) = \mathrm{O}\) , which indeed proves that \(\overline{\mathrm{O}} -\mathrm{O}\subset \mathrm{X}\) .

Let X be a compact subset of C. Let \(O_{\infty}\) be the unbounded connected component of \(\mathbf{C} - \mathbf{X}\) , and let \((\mathrm{O}_{i})_{i\in\mathrm{I}}\) be the family of bounded connected components of \(\mathbf{C} - \mathbf{X}\) , the sets \(O_{i}\) being pairwise distinct.

Let E be a subset of \(\mathbf{C} - \mathbf{X}\) . We denote by \(\mathrm{R}_{\mathrm{E}}(\mathrm{X})\) the closure in \(\mathcal{C}(\mathrm{X})\) of the set of functions \(f|\mathrm{X}\) , where \(f\) is a rational function on \(\mathbf{C}\) all of whose poles belong to E. The algebra \(\mathrm{R}_{\mathrm{E}}(\mathrm{X})\) is a unital Banach subalgebra of \(\mathcal{C}(\mathrm{X})\) which separates the points of X. Let z be the identity function on X. The full closed subalgebra of \(\mathrm{R}_{\mathrm{E}}(\mathrm{X})\) generated by z is equal to \(\mathrm{R}_{\mathrm{E}}(\mathrm{X})\) (lemma 2 of I, p.~6). The elements of \(\mathrm{R}_{\mathrm{E}}(\mathrm{X})\) are holomorphic in the interior of X.

In particular, we have \(\mathrm{R}_{\varnothing}(\mathrm{X})=\mathrm{P}(\mathrm{X})\) . We denote \(\mathrm{R}(\mathrm{X})=\mathrm{R}_{\mathbf{C}-\mathbf{X}}(\mathrm{X})\) . We denote by I(E) the set of elements \(i\in I\) such that \(E\cap O_{i}=\varnothing\) , and

\[
\mathrm{X} _ {\mathrm{E}} = \mathrm{X} \cup \Big (\bigcup_ {i \in \mathrm{I} (\mathrm{E})} \mathrm{O} _ {i} \Big).
\]

The set \(X_{E}\) is compact, since it is bounded and closed, its complement in \(\mathbf{C}\) being the union of the open \(O_{\infty}\) and open sets \(O_{i}\) that meet E.

PROPOSITION 5. --- With the above notation:

\begin{enumerate}
\def\labelenumi{\alph{enumi})}
\item
  The restriction map h of \(\mathrm{R}_{\mathrm{E}}(\mathrm{X}_{\mathrm{E}})\) into \(\mathrm{R}_{\mathrm{E}}(\mathrm{X})\) is an isometric isomorphism;
\item
  The algebra \(\mathrm{R}_{\mathrm{E}}(\mathrm{X}_{\mathrm{E}})\) is a full subalgebra of \(\mathcal{C}(\mathrm{X}_{\mathrm{E}})\) ;
\item
  Every character of \(\mathrm{R}_{\mathrm{E}}(\mathrm{X}_{\mathrm{E}})\) is of the form \(f \mapsto f(w)\) , where w is an element of \(X_{E}\) ;
\item
  The map \(\chi \mapsto \chi (z)\) is a homeomorphism from \(\mathsf{X}(\mathrm{R}_{\mathrm{E}}(\mathrm{X}))\) onto \(\mathrm{X}_{\mathrm{E}}\) ;
\item
  If E' is a subset of \(\mathbf{C} - \mathbf{X}\) , then we have \(\mathrm{R}_{\mathrm{E}}(\mathrm{X}) = \mathrm{R}_{\mathrm{E}'}(\mathrm{X})\) if and only if \(\mathrm{X}_{\mathrm{E}} = \mathrm{X}_{\mathrm{E}'}\) , which is also equivalent to \(\mathrm{I}(\mathrm{E}) = \mathrm{I}(\mathrm{E}')\) .
\end{enumerate}

The restriction map \(h\) of \(\mathrm{R}_{\mathrm{E}}(\mathrm{X}_{\mathrm{E}})\) into \(\mathrm{R}_{\mathrm{E}}(\mathrm{X})\) is a Banach algebra morphism such that \(\| h(f)\| \leqslant \| f\|\) for every function \(f\in \mathrm{R}_{\mathrm{E}}(\mathrm{X}_{\mathrm{E}})\) . Let \(f\in \mathrm{R}_{\mathrm{E}}(\mathrm{X}_{\mathrm{E}})\) . Let \(i\in \mathrm{I}(\mathrm{E})\) . Since \(f\) is holomorphic in a neighborhood of the bounded open set \(\mathrm{O}_i\) , the maximum principle (VAR, I, p.~29) implies that there exists an element \(z_0\) in the boundary of \(\mathrm{O}_i\) such that \(|f(z_0)| = \sup_{z\in \mathrm{O}_i}|f(z)|\) . Since the boundary of \(\mathrm{O}_i\) is contained in X (lemma 1), it follows that \(\sup_{z\in \mathrm{O}_i}|f(z)|\leqslant \| h(f)\|\) . Since this inequality holds for every \(i\in \mathrm{I}(\mathrm{E})\) , it results that \(\| f\| \leqslant \| h(f)\|\) . The morphism \(h\) is therefore isometric, and in particular injective. Let us now prove that it is surjective. Let \(g\in \mathrm{R}_{\mathrm{E}}(\mathrm{X})\) . There exists a sequence \((f_n)\) of rational functions whose poles belong to E which converges uniformly to \(g\) on X. The \(f_n|\mathrm{X}_{\mathrm{E}}\) are elements of \(\mathrm{R}_{\mathrm{E}}(\mathrm{X}_{\mathrm{E}})\) . Since \(h\) is isometric, the sequence \((f_n|\mathrm{X}_{\mathrm{E}})\) converges in \(\mathcal{C}(\mathrm{X}_{\mathrm{E}})\) . If \(f\) is its limit, we have \(f\in \mathrm{R}_{\mathrm{E}}(\mathrm{X}_{\mathrm{E}})\) and \(g = f|\mathrm{X} = h(f)\) . Hence \(a\) ) .

Let us prove assertion d). Apply prop. 3 of I, p.~144 to the algebra \(B = \mathrm{R}_{\mathrm{E}}(\mathrm{X})\) . Assertion b) of loc. cit. implies that the map which associates \(\chi(z)\) to \(\chi\) is a homeomorphism from \(\mathsf{X}(\mathrm{R}_{\mathrm{E}}(\mathrm{X}))\) onto \(\mathrm{Sp}_{\mathrm{R}_{\mathrm{E}}(\mathrm{X})}(z)\) . Let \(z_{\mathrm{E}}\) be the identity map of \(\mathrm{X}_{\mathrm{E}}\) . By loc. cit., a), we have \(\mathrm{Sp}_{\mathrm{R}_{\mathrm{E}}(\mathrm{X})}(z) = \mathrm{Sp}_{\mathrm{R}_{\mathrm{E}}(\mathrm{X}_{\mathrm{E}})}(z_{\mathrm{E}})\) . It therefore suffices to show that \(\mathrm{Sp}_{\mathrm{R}_{\mathrm{E}}(\mathrm{X}_{\mathrm{E}})}(z_{\mathrm{E}}) = \mathrm{X}_{\mathrm{E}}\) . The corollary of prop. 6 of I, p.~28 shows that \(\mathrm{Sp}_{\mathrm{R}_{\mathrm{E}}(\mathrm{X}_{\mathrm{E}})}(z_{\mathrm{E}})\) is the union of \(\mathrm{Sp}_{\mathcal{C}(\mathrm{X}_{\mathrm{E}})}(z_{\mathrm{E}}) = \mathrm{X}_{\mathrm{E}}\) and of certain bounded connected components of the complement of \(\mathrm{X}_{\mathrm{E}}\) . Let \(\mathrm{O}_{i}\) be one of the bounded connected components of the complement of \(\mathrm{X}_{\mathrm{E}}\) . The intersection \(\mathrm{E} \cap \mathrm{O}_{i}\) is therefore nonempty. Let \(\lambda \in \mathrm{E} \cap \mathrm{O}_{i}\) ; since \((\lambda - z_{\mathrm{E}})^{-1} \in \mathrm{R}_{\mathrm{E}}(\mathrm{X}_{\mathrm{E}})\) , we have \(\lambda \notin \mathrm{Sp}_{\mathrm{R}_{\mathrm{E}}(\mathrm{X}_{\mathrm{E}})}(z_{\mathrm{E}})\) . Thus, \(\mathrm{O}_{i}\) is not contained in \(\mathrm{Sp}_{\mathrm{R}_{\mathrm{E}}(\mathrm{X}_{\mathrm{E}})}(z_{\mathrm{E}})\) . This shows that \(\mathrm{Sp}_{\mathrm{R}_{\mathrm{E}}(\mathrm{X}_{\mathrm{E}})}(z_{\mathrm{E}}) = \mathrm{X}_{\mathrm{E}}\) .

This equality also implies that \(\mathrm{Sp}_{\mathrm{R}_{\mathrm{E}}(\mathrm{X}_{\mathrm{E}})}(z_{\mathrm{E}})=\mathrm{Sp}_{\mathcal{C}(\mathrm{X}_{\mathrm{E}})}(z_{\mathrm{E}})\) . This establishes condition (iii) of prop. 11 of I, p.~41, applied to the canonical injection of \(\mathrm{R}_{\mathrm{E}}(\mathrm{X}_{\mathrm{E}})\) into \(\mathcal{C}(\mathrm{X}_{\mathrm{E}})\) and to the element \(z_{\mathrm{E}}\) . Assertions b) and c) are the equivalent conditions (i) and (ii) of loc. cit.

Assertion \(d\) ) shows that \(\mathrm{X}_{\mathrm{E}} = \mathrm{X}_{\mathrm{E}'}\) if \(\mathrm{R}_{\mathrm{E}}(\mathrm{X}) = \mathrm{R}_{\mathrm{E}'}(\mathrm{X})\) . Conversely, suppose that \(\mathrm{X}_{\mathrm{E}} = \mathrm{X}_{\mathrm{E}'}\) . By replacing \(\mathrm{E}'\) by \(\mathrm{E} \cup \mathrm{E}'\) , we may assume that \(\mathrm{E} \subset \mathrm{E}'\) . According to \(b\) ), the algebra \(\mathrm{R}_{\mathrm{E}}(\mathrm{X}_{\mathrm{E}})\) is a full closed subalgebra of \(\mathcal{C}(\mathrm{X}_{\mathrm{E}})\) , and therefore also of \(\mathrm{R}_{\mathrm{E}'}(\mathrm{X}_{\mathrm{E}})\) . It contains \(z_{\mathrm{E}}\) , and hence \(\mathrm{R}_{\mathrm{E}}(\mathrm{X}_{\mathrm{E}}) = \mathrm{R}_{\mathrm{E}'}(\mathrm{X}_{\mathrm{E}})\) . Applying \(a\) ), we deduce that \(\mathrm{R}_{\mathrm{E}}(\mathrm{X}) = \mathrm{R}_{\mathrm{E}'}(\mathrm{X})\) . Finally, the equivalence of \(\mathrm{X}_{\mathrm{E}} = \mathrm{X}_{\mathrm{E}'}\) and \(\mathrm{I}(\mathrm{E}) = \mathrm{I}(\mathrm{E}')\) is a consequence of the definitions.

COROLLARY 1. --- The following conditions are equivalent:

\begin{enumerate}
\def\labelenumi{(\roman{enumi})}
\item
  The set E meets every bounded connected component of \(\mathbf{C} - \mathbf{X}\) ;
\item
  The map \(\chi \mapsto \chi (z)\) is a homeomorphism from \(\mathsf{X}(\mathrm{R}_{\mathrm{E}}(\mathrm{X}))\) onto \(\mathrm{X}\) ;
\item
  We have \(\mathrm{R_E(X)} = \mathrm{R(X)}\) .
\end{enumerate}

Let \(E' = \mathbf{C} - \mathbf{X}\) . Conditions (i), (ii), and (iii) are respectively equivalent to \(\mathrm{I}(\mathrm{E}) = \mathrm{I}(\mathrm{E}')\) , to \(\mathrm{X}_{\mathrm{E}} = \mathrm{X}_{\mathrm{E}'}\) (by prop. 5, d)) and to \(\mathrm{R}_{\mathrm{E}}(\mathrm{X}) = \mathrm{R}_{\mathrm{E}'}(\mathrm{X})\) . They are therefore equivalent to each other by prop. 5, e).

The following corollary makes the Runge theorem more precise (th. 3 of I, p.~69).

COROLLARY 2 (Runge's theorem). --- For every \(i \in I\) , let \(\lambda_{i}\) be a point of \(O_{i}\) . Let f be a complex function holomorphic in an open neighborhood of X. Then \(f|X\) is the uniform limit on X of restrictions of rational functions whose poles are among the \(\lambda_{i}\) .

With \(E = \{\lambda_{i}\}_{i \in I}\) , the hypothesis is condition (i) of Cor. 1. We therefore have \(\mathrm{R}_{\mathrm{E}}(\mathrm{X}) = \mathrm{R}(\mathrm{X})\) , and Th. 3 of I, p.~69 shows that \(\mathrm{R}(\mathrm{X}) = \mathcal{O}(\mathrm{X})\) .

